To begin we need to add the OpenCHAMI RPM to the SMS node, by adding it to the list of
available packages locally. This requires network access from the {\em master}
node to the internet.

% begin_ohpc_run
% ohpc_validation_newline
% ohpc_comment_header Enable OpenCHAMI Repositories and Support \ref{sec:enable_openchami}
\begin{lstlisting}[language=bash,keywords={},basicstyle=\fontencoding{T1}\fontsize{8.0}{10}\ttfamily,
	literate={VER}{\OHPCVerTree{}}1 {OSREPO}{\OSTree{}}1 {TAG}{\OSTag{}}1 {ARCH}{\arch{}}1 {-}{-}1]i
[sms](*\#*) (*\install*) podman buildah ansible-core nfs-utils tcpdump tftp curl yq
[sms](*\#*) rpm --import https://github.com/OpenCHAMI/release/releases/download/vVERLONG/public_gpg_key.asc
[sms](*\#*) (*\install*) https://github.com/OpenCHAMI/release/releases/download/vVERLONG/openchami-VERLONG.rpm
\end{lstlisting}
% ohpc_validation_newline
% end_ohpc_run
